% Signal Hierarchy Theory

\subsection{Motivation and Intuition}

Cellular regulation exhibits hierarchical organization where higher-layer signals control lower-layer processes through information channels with consumption semantics. Existing Bio-PN formalisms separate metabolic mass transfer (normal arcs) from regulatory catalysis (test arcs), but cannot model regulatory consumption---the depletion of regulatory molecules during decision-making that ensures commitment finality. This limitation prevents unified modeling where metabolic state (e.g., ATP concentration) functions simultaneously as metabolite and hierarchical control signal. Bacteriophage lambda lysis-lysogeny provides an illustrative external example: phage-encoded CI and Cro regulatory proteins function as signals whose consumption determines bistable commitment between developmental pathways. However, \textit{Bacillus subtilis} sporulation best exemplifies this hierarchical architecture: ATP availability functions as both metabolic resource and regulatory signal controlling a five-layer hierarchical decision between competence and sporulation pathways, with distinct attractor basins ensuring commitment finality.

\begin{center}
\small
\begin{tabular}{@{}ll@{}}
\multicolumn{2}{c}{\textbf{\textit{B. subtilis} Sporulation Hierarchy}} \\
\midrule
\textbf{L4:} & Metabolic state (ATP) \\
\textbf{L3:} & Energy sensing \\
\textbf{L2:} & Phosphorelay integration \\
\textbf{L1:} & Sigma factor decision \\
\textbf{L0:} & Pathway execution
\end{tabular}
\end{center}

\vspace{3pt}
\noindent Critical distinction: ATP signal is \textit{consumed} during decision-making, creating distinct attractor basins. This five-layer architecture demonstrates the formalism's capability to model complex hierarchical regulatory control in bacterial developmental decisions. This differs from catalytic test arcs (non-consuming enzyme reads). Signal consumption enables hierarchical preemption---higher layers collapse decision spaces at lower layers.

\subsection{Signal Places: Formal Definition}

\begin{definition}[Signal Place]
A place $p \in P$ is a \textit{signal place} ($p \in \Psi \subseteq P$) if:
\textbf{(1) $p$ appears in rate function $\Phi(t)$ of some transition $t \in T$}--- \textbf{(2) $p$ is \textbf{not} connected to $t$ by normal arc}--- \textbf{(3) $p$ represents information channel, not mass transfer}---
\end{definition}

\vspace{3pt}
\noindent We identify four signal types in biological systems:

\vspace{3pt}
\noindent\textbf{(1) \textbf{Energy signals}}---ATP, NADH, energy charge. Hierarchically regulate anabolic/catabolic flux. \textbf{(2) \textbf{Spatial signals}}---Compartment markers, membrane integrity, diffusion gradients. Control localization-dependent processes. \textbf{(3) \textbf{Quorum signals}}---Autoinducers, population density markers. Enable coordinated multicellular behavior. \textbf{(4) \textbf{Regulatory signals}}---Transcription factors, checkpoint proteins. Implement decision logic and pathway switching.

\vspace{3pt}
\noindent\textbf{Mathematical integration}: Rate functions extend to include signal dependencies:
\begin{equation}
\label{eq:rate_signal}
\Phi(t) = k_t \cdot f({}^\bullet t) \cdot g(\Sigma(t)) \cdot h(\Psi(t))
\end{equation}
where $f$ models substrate kinetics, $g$ models regulatory modulation, and $h$ models signal influence.

\subsection{Signal Flow Arcs: Semantics}

\begin{definition}[Signal Flow Arc]
An arc $a \in F$ has type \textit{signal\_flow} if:
\textbf{Source or target is signal place ($p \in \Psi$)}--- \textbf{Consumes tokens during transition firing (unlike test arcs)}--- \textbf{Represents information transfer with depletion}---
\end{definition}

\vspace{3pt}
\noindent\textbf{Execution semantics} (two-phase):

\noindent\textbf{Phase 1 (Enabling)}---Check signal availability
\begin{equation}
\label{eq:signal_enable}
M(p_{\text{signal}}) \geq E(t, p_{\text{signal}})
\end{equation}
where $E: T \times \Psi \to \mathbb{R}_{\geq 0}$ defines enablement thresholds.

\noindent\textbf{Phase 2 (Firing)}---Consume signal tokens
\begin{equation}
\label{eq:signal_consume}
M'(p_{\text{signal}}) = M(p_{\text{signal}}) - W(p_{\text{signal}}, t)
\end{equation}
where $W$ defines arc weight (consumption amount).

\subsection{Comparison with Test Arcs}

Test arcs in classical Bio-PN model non-consuming reads (catalysis), while signal flow arcs model consuming information transfer. This distinction is critical for hierarchical control and decision modeling:

\begin{table}[h]
\caption{Test vs signal flow arcs}
\label{tab:test_vs_signal}
\centering
\small
\begin{tabular}{@{}p{0.35\columnwidth}p{0.28\columnwidth}p{0.28\columnwidth}@{}}
\toprule
\textbf{Property} & \textbf{Test Arc} & \textbf{Signal Flow Arc} \\
\midrule
Token consumption & Non-consuming & Consuming \\
Biological semantics & Catalytic effect & Information transfer \\
Primary use case & Enzyme catalysis & Signal consumption \\
Hierarchical control & Cannot preempt & Enables preemption \\
Decision modeling & Observation only & Exclusive pathways \\
\bottomrule
\end{tabular}
\end{table}

\vspace{3pt}
\subsection{Hierarchical Layers}

\begin{definition}[Hierarchical Layer Assignment]
Function $L: \Psi \to \mathbb{N}$ assigns layer indices to signal places where:
\begin{equation}
\label{eq:layer_assignment}
L(p_1) < L(p_2) \iff p_1 \text{ controls } p_2
\end{equation}
via signal flow paths $p_1 \to \cdots \to p_2$.
\end{definition}

\vspace{3pt}
\noindent\textbf{Layer computation algorithm}: Given annotated signal places $\Psi$ and transitions $T$:

\vspace{2pt}
\noindent\textbf{(1)} Construct signal dependency graph $G = (\Psi, E)$ where edge $(p_i, p_j) \in E$ exists if signal $p_i$ controls (via enablement threshold or rate modulation) transitions producing/consuming signal $p_j$

\vspace{2pt}
\noindent\textbf{(2)} Verify acyclicity: If $G$ contains cycle, hierarchy undefined (feedback loops require fixpoint semantics---2\% of analyzed models)

\vspace{2pt}
\noindent\textbf{(3)} Topological sort: Order signal places by dependencies (Kahn's algorithm)

\vspace{2pt}
\noindent\textbf{(4)} Assign layers: $L(p) = $ longest path from $p$ to terminal signals. Terminal signals (sinks) receive layer 0

\vspace{3pt}
\noindent Complexity: $O(|\Psi| + |E|)$ linear in signal network size.

\begin{theorem}[Layer Assignment Consistency]\footnote{Complete proofs establishing soundness and completeness of the layer assignment algorithm, including treatment of edge cases and complexity analysis, are provided in~\cite{subtilis2026}.}
For acyclic signal dependency graph $G = (\Psi, E)$, the layer assignment $L: \Psi \to \mathbb{N}$ exists, is unique, and satisfies the hierarchical constraint: $\forall (p_i, p_j) \in E: L(p_i) > L(p_j)$.
\end{theorem}

\begin{proof}[Proof sketch]
Acyclicity guarantees topological ordering exists. Longest-path construction ensures uniqueness: any signal $p$ receives layer equal to maximum layer of signals it controls plus one. Hierarchical constraint satisfied by construction since controlling signals receive strictly higher layers than controlled signals.
\end{proof}

\begin{figure}[t]
\centering
\begin{tikzpicture}[scale=0.7, transform shape,
  node distance=1.5cm,
  place/.style={circle, draw=black, thick, minimum size=8mm},
  signal/.style={regular polygon, regular polygon sides=6, draw=blue!70, thick, minimum size=8mm, fill=blue!10},
  transition/.style={rectangle, draw=black, thick, minimum width=5mm, minimum height=10mm, fill=black},
  arc/.style={->, thick},
  signal_arc/.style={->, thick, blue!70, dashed}
]

% Layer 2 (Energy Sensing)
\node[signal] (atp_signal) at (0.5, 4) {$ATP_s$};
\node at (4, 4) {\small\textbf{Layer 2}};

% Layer 1 (Decision)
\node[transition] (anabolic) at (-1.5, 2.5) {};
\node[left=1mm of anabolic] {\tiny Anabolic};
\node[transition] (catabolic) at (1.5, 2.5) {};
\node[right=1mm of catabolic] {\tiny Catabolic};
\node at (4, 2.5) {\small\textbf{Layer 1}};

% Layer 0 (Execution)
\node[place] (biosyn) at (-1.5, 0.5) {$B$};
\node[below=1mm of biosyn] {\tiny Biosynthesis};
\node[place] (degr) at (1.5, 0.5) {$D$};
\node[below=1mm of degr] {\tiny Degradation};
\node at (4, 0.5) {\small\textbf{Layer 0}};

% Metabolic layer (showing ATP as metabolite)
\node[place] (glc) at (-3, 4) {$Glc$};
\node[transition] (glyc) at (-1.5, 4) {};

% Signal flow arcs (blue dashed)
\draw[signal_arc] (atp_signal) -- node[midway, right, font=\tiny, yshift=-6pt] {$E=2.5$} (anabolic);

% Normal arcs (black solid)
\draw[arc] (anabolic) -- (biosyn);
\draw[arc] (biosyn) -- (catabolic);
\draw[arc] (catabolic) -- (degr);
\draw[arc] (glc) -- (glyc);
\draw[arc] (glyc) -- (atp_signal);
\draw[arc] (biosyn) to[out=180, in=180] (glc);
\draw[arc] (degr) -- (atp_signal);

\end{tikzpicture}
\caption{\textbf{Energy-gated metabolic switch (illustrative model).} ATP signal place ($ATP_s$, hexagon) serves dual role: receives tokens from glycolysis (glucose metabolism) and catabolic degradation (biomass breakdown), while controlling pathway selection via signal flow arc (blue dashed). High ATP ($\geq 2.5$ mM) enables anabolic biosynthesis through signal consumption; ATP depletion below threshold blocks anabolic transition, allowing constitutive catabolism to degrade biosynthesis products and regenerate ATP. Signal consumption enables hierarchical preemption---test arcs (non-consuming) cannot model exclusive pathway commitment. Layer assignment: Energy Sensing (Layer 2) $>$ Decision (Layer 1) $>$ Execution (Layer 0).}
\label{fig:shypn_example}
\end{figure}

\subsection{Hierarchical Preemption Mechanism}

Execution order respects layer hierarchy:

\noindent\textbf{(1) Priority ordering}---Higher layers (larger $L$) fire first

\vspace{4pt}
\noindent\textbf{(2) Signal consumption}---Consumption at layer $k$ blocks dependent transitions at layer $k-1$

\vspace{4pt}
\noindent\textbf{(3) Decision collapse}---Signal availability at higher layer preempts lower-layer alternatives

\vspace{4pt}
\begin{theorem}[Preemption Correctness]\footnote{Detailed proof including treatment of concurrent transition firing, tie-breaking rules, and stochastic extensions is available in~\cite{subtilis2026}.}
If transition $t$ at layer $k$ consumes signal $p_s$ enabling transition $t'$ at layer $k' < k$, then $t'$ cannot fire after $t$ fires and depletes $p_s$ below threshold $E(t', p_s)$.
\end{theorem}

\begin{proof}[Proof sketch]
Transition $t'$ enabled requires $M(p_s) \geq E(t', p_s)$. If $t$ fires consuming $W(p_s, t)$ tokens, resulting marking $M'(p_s) = M(p_s) - W(p_s, t)$. If $M'(p_s) < E(t', p_s)$, then $t'$ disabled. This models hierarchical preemption: higher-layer signal consumption blocks lower-layer alternatives.
\end{proof}

\vspace{3pt}
\noindent\textbf{Example}: In \textit{B. subtilis} sporulation (detailed in Section 5), ATP signal at layer 2 gates phosphorelay activation at layer 1. ATP depletion below threshold preempts sporulation pathway, activating competence alternative. Signal consumption formalizes energy-controlled developmental commitment.

\vspace{3pt}
\noindent This hierarchical preemption cannot be modeled with test arcs (non-consuming). Signal consumption is essential for decision finality. Figure~\ref{fig:shypn_example} illustrates this formalism through an energy-gated metabolic switch showing the interplay between signal places, signal flow arcs, and hierarchical layers.
